\hypertarget{abhaken}{}\einheitenkopf*{Das kann ich}
\zweigzeile{Hake am Ende ab, was du kannst.}
{\small
\par\medskip\noindent\textbf{Kennst du schon}\par\smallskip
\abhak{1}{Ich kann Dezimalzahlen malnehmen und das Ergebnis sinnvoll runden.}
\abhak{2}{Ich kann einen Wert um einen Prozentsatz erhöhen oder senken, indem ich mit einer Dezimalzahl malnehme.}
\abhak{3}{Ich finde den Fehler beim Erhöhen um Prozent.}
\abhak{4}{Ich kann Punkte in einem Koordinatensystem mit großen Schritten ablesen.}
\abhak{5}{Ich kann eine Reihe fortsetzen, in der jedes Mal derselbe Betrag dazukommt oder wegfällt.}
\abhak{6}{Ich kann zwei Werte durcheinander teilen und das Ergebnis deuten.}
\abhak{7}{Ich kann Prozentwert und Prozentsatz berechnen.}
\abhak{8}{Ich kann ausrechnen, wie ein Guthaben mit Zinseszins wächst.}
\abhak{9}{Ich kann Potenzen mit dem Taschenrechner berechnen.}
\par\medskip\noindent\textbf{Einheit 1 von 5 · Lineares und exponentielles Wachstum unterscheiden}\par\smallskip
\abhak{10}{Ich erkenne, ob sich ein Wert jedes Mal um denselben Betrag oder um denselben Prozentsatz ändert.}
\abhak{11}{Ich erkenne, ob in einer Reihe immer dasselbe dazukommt oder immer mit derselben Zahl malgenommen wird.}
\abhak{12}{Ich erkenne, wo ein Graph die senkrechte Achse trifft.}
\abhak{13}{Ich kann an einer Tabelle entscheiden, ob ein Wert linear oder exponentiell wächst.}
\abhak{14}{Ich kann einen Graphen zu einem Wachstum auswählen.}
\abhak{15}{Ich kann einen Graphen zu einem Wachstum auswählen – weiter.}
\abhak{16}{Ich kann Wertepaare eines Wachstums als Punkte in ein Koordinatensystem eintragen.}
\abhak{17}{Ich kann Wertepaare eines Wachstums als Punkte in ein Koordinatensystem eintragen – weiter.}
\abhak{18}{Ich finde den Fehler bei der Entscheidung zwischen linear und exponentiell.}
\abhak{19}{Ich kann begründen, warum gleicher Prozentsatz nicht gleicher Betrag heißt.}
\abhak{20}{Ich kann mit linearem und exponentiellem Wachstum eine Entscheidung treffen.}
\par\medskip\noindent\textbf{Einheit 2 von 5 · Wachstumsfaktor und Wachstumstabelle}\par\smallskip
\abhak{21}{Ich erkenne am Faktor, ob ein Wert zunimmt oder abnimmt.}
\abhak{22}{Ich kann den Wachstumsfaktor bestimmen.}
\abhak{23}{Ich kann eine Wachstumstabelle mit dem Wachstumsfaktor fortschreiben.}
\abhak{24}{Ich kann eine Wachstumstabelle mit dem Wachstumsfaktor fortschreiben – weiter.}
\abhak{25}{Ich finde den Fehler beim Fortschreiben einer Wachstumstabelle.}
\abhak{26}{Ich kann begründen, warum der Quotient den Wachstumsfaktor zeigt.}
\abhak{27}{Ich kann mit einer Wachstumstabelle eine Entscheidung treffen.}
\par\medskip\noindent\textbf{Einheit 3 von 5 · Exponentialfunktion aufstellen und auswerten}\par\smallskip
\abhak{28}{Ich erkenne, welches Jahr der Schritt null ist.}
\abhak{29}{Ich kann eine Exponentialfunktion aufstellen.}
\abhak{30}{Ich kann eine Exponentialfunktion aufstellen – weiter.}
\abhak{31}{Ich kann mit einer Exponentialfunktion Werte berechnen.}
\abhak{32}{Ich kann mit einer Exponentialfunktion Werte berechnen – weiter.}
\abhak{33}{Ich finde den Fehler beim Berechnen mit einer Exponentialfunktion.}
\abhak{34}{Ich kann begründen, warum bei exponentiellem Wachstum eine Potenz entsteht.}
\abhak{35}{Ich kann mit einer Exponentialfunktion eine Entscheidung treffen.}
\par\medskip\noindent\textbf{Einheit 4 von 5 · Verdopplungs- und Halbwertszeit}\par\smallskip
\abhak{36}{Ich erkenne, ob in einer Frage die Zeit oder der Bestand verdoppelt oder halbiert wird.}
\abhak{37}{Ich kann die Verdopplungszeit und die Halbwertszeit bestimmen.}
\abhak{38}{Ich kann die Verdopplungszeit und die Halbwertszeit bestimmen – weiter.}
\abhak{39}{Ich kann die Verdopplungszeit und die Halbwertszeit bestimmen – weiter.}
\abhak{40}{Ich finde den Fehler beim Bestimmen der Halbwertszeit.}
\abhak{41}{Ich kann begründen, warum die Verdopplungszeit immer gleich lang ist.}
\abhak{42}{Ich kann mit der Halbwertszeit eine Entscheidung treffen.}
\par\medskip\noindent\textbf{Einheit 5 von 5 · Potenzfunktionen mit natürlichem Exponenten}\par\smallskip
\abhak{43}{Ich kann Wertetabellen zu Potenzfunktionen ausfüllen.}
\abhak{44}{Ich kann Funktionswerte von Potenzfunktionen berechnen.}
\abhak{45}{Ich kann den Graphen einer Potenzfunktion zeichnen.}
\abhak{46}{Ich kann an der Gleichung erkennen, wie der Graph einer Potenzfunktion verläuft.}
\abhak{47}{Ich kann Gleichungen und Graphen von Potenzfunktionen einander zuordnen.}
\abhak{48}{Ich kann unterscheiden, ob die Variable in der Basis oder im Exponenten steht.}
\abhak{49}{Ich finde den Fehler beim Berechnen von Funktionswerten.}
\abhak{50}{Ich kann Eigenschaften von Potenzfunktionen begründen.}
\abhak{51}{Ich kann mit einer Potenzfunktion eine Sachaufgabe lösen.}
}
